\section*{Chapitre \ref{ch:iv1:conditional-probability-and-bayes} --- Probabilité conditionnelle et Bayes}
\markboth{\protect\omTocMark{}{Réponses~: Probabilité conditionnelle et Bayes}}{\protect\omTocMark{}{Réponses~: Probabilité conditionnelle et Bayes}}

\begin{solution}{iq:iv1:conditional-probability-and-bayes:1}
Si l'on sait que «~au moins un est une fille~», la probabilité vaut $1/3$. Si l'on a rencontré un enfant pris au hasard qui est une fille, elle vaut $1/2$.

Les quatre familles BB, BG, GB, GG sont équiprobables. «~Au moins une fille~» élimine seulement BB, et parmi les trois restantes, une seule est GG. Lorsqu'on rencontre un enfant pris au hasard, on en apprend davantage. La famille GG montre une fille à coup sûr, alors que BG et GB n'en montrent une que dans la moitié des cas~: les poids valent donc $1 : \tfrac12 : \tfrac12$, et GG a une probabilité de $1/2$.

\iqwrong{Supposer que les deux énoncés sont le même fait. Ils diffèrent par la manière dont l'information est parvenue\index{protocole d'information}~: ce qui compte est la probabilité d'entendre ce message sous chacune des hypothèses, et non seulement sa véracité.}
\end{solution}

\enlargethispage{-3\baselineskip}%
\begin{solution}{iq:iv1:conditional-probability-and-bayes:2}
La probabilité vaut $1/6$, soit environ 17~\%.

Utiliser Bayes par les cotes\index{Bayes par les cotes}. La cote a priori de la maladie vaut $1 : 99$. Un test positif est $0.99/0.05 = 19.8$ fois plus probable avec la maladie que sans elle, ce qui est donc le rapport de vraisemblance\index{rapport de vraisemblance}. La cote a posteriori vaut $19.8 : 99 = 1 : 5$, soit une probabilité de $1/6$.

Concrètement~: sur 10\,000 personnes, 100 sont malades et 99 d'entre elles testent positives. Parmi les 9\,900 bien portantes, 495 testent positives. Donc 99 des 594 positifs sont malades.

\iqwrong{Répondre 99~\%, ou un chiffre proche. Cela confond $P(\text{positif} \mid \text{malade})$ avec $P(\text{malade} \mid \text{positif})$, et ignore la rareté de la maladie.}
\iqwhy{La même somme fixe la valeur de toute alerte. Un signal juste dans 99~\% des cas pour des événements rares produit tout de même surtout de fausses alarmes.}
\end{solution}

\begin{solution}{iq:iv1:conditional-probability-and-bayes:3}
Changer de porte~: on gagne avec une probabilité de $2/3$. Avec 100 portes, on gagne avec une probabilité de $99/100$.

Le premier choix est le bon avec une probabilité de $1/3$. L'action de l'animateur ne dit rien de la porte choisie, car il peut toujours ouvrir une porte derrière laquelle il n'y a rien, quelle que soit celle que l'on a retenue. Cette porte conserve donc son $1/3$, et la porte restante recueille les $2/3$ restants. De manière équivalente, changer de porte gagne exactement lorsque le premier choix était mauvais.

La version à cent portes rend cela évident~: la porte choisie est la bonne avec une probabilité de $1/100$, et l'animateur, qui sait, a montré où se trouverait le lot dans tous les autres cas.

\iqwrong{«~Il reste deux portes, donc chacune a une chance sur deux.~» Deux issues ne signifient pas des chances égales. Le savoir de l'animateur est ce qui rend la mise à jour inégale, et un animateur qui ouvrirait au hasard donnerait une réponse différente.}
\end{solution}

\begin{solution}{iq:iv1:conditional-probability-and-bayes:4}
La probabilité vaut $1024/1123 \approx 0.912$.

Par les cotes\index{Bayes par les cotes}, la cote a priori de la pièce à deux faces vaut $1 : 99$. Chaque pile est deux fois plus probable avec cette pièce, donc chaque pile double la cote. Après 10 piles, la cote vaut $1024 : 99$, ce qui donne $1024/1123$.

La croyance dépasse un demi entre la sixième et la septième pile, puisque $2^6 = 64 < 99 < 128 = 2^7$.

\iqwhy{Les cotes font de cette mise à jour un simple calcul et montrent à quelle vitesse une preuve répétée écrase un a priori. Dix doublements multiplient par environ mille, si bien qu'un a priori de 1 sur 100 tombe vite, mais qu'un a priori de 1 sur un million ne tomberait pas aussi vite.}
\end{solution}

\begin{solution}{iq:iv1:conditional-probability-and-bayes:5}
La probabilité vaut $2/3$.

Compter les faces rouges, et non les cartes. Il y a trois faces rouges, et chacune a autant de chances d'être celle que l'on regarde. Deux de ces trois faces appartiennent à la carte rouge-rouge, et leur autre face est rouge. Une seule appartient à la carte mixte, et son autre face est bleue.

\iqwrong{Une moitié, car seules deux cartes pourraient être posées sur la table. Cela traite les deux cartes possibles comme équiprobables, alors que la carte rouge-rouge a deux fois plus de chances de montrer une face rouge.}

\end{solution}

\begin{solution}{iq:iv1:conditional-probability-and-bayes:6}
La probabilité vaut $12/29 \approx 0.414$, si bien que le taxi reste plus probablement vert.

Utiliser les cotes\index{Bayes par les cotes}. La cote a priori du bleu vaut $15 : 85$. Le témoin a 4 fois plus de chances de dire bleu quand le taxi est bleu que quand il est vert ($0.8/0.2$). Donc la cote a posteriori vaut $60 : 85 = 12 : 17$, ce qui donne $12/29$.

\iqwrong{80~\%, la fiabilité du témoin. Cela répond à «~à quelle fréquence le témoin a raison au sujet d'un taxi quelconque~», et non à «~étant donné cette déclaration sur une couleur rare, à quelle fréquence a-t-elle raison~». Les taxis verts sont si nombreux que leurs erreurs dépassent les bonnes déclarations sur les bleus, 17 contre 12.}
\iqfollow{Et si le témoin avait raison dans la moitié des cas~?}{Alors le rapport de vraisemblance vaut 1, et la réponse reste à l'a priori, soit 15~\%.}
\end{solution}

\begin{solution}{iq:iv1:conditional-probability-and-bayes:7}
La probabilité d'obtenir deux six vaut $1/11$, et celle d'une somme égale à 7 vaut $2/11$.

Sur les 36 lancers, 11 contiennent un six~: 6 avec le premier dé à six, 6 avec le second dé à six, la configuration $(6, 6)$ étant comptée deux fois. Une seule des 11 est $(6, 6)$, et deux d'entre elles font une somme de 7, à savoir $(1, 6)$ et $(6, 1)$.

\iqwrong{Un sixième, en raisonnant ainsi~: «~un dé est un six, donc l'autre est un dé indépendant~». Cela n'est juste que si l'on sait que le premier dé est un six. «~Au moins un~» ne porte pas sur un dé précis, et cela modifie le décompte.}
\iqfollow{Et si la personne regarde au hasard un dé et affirme qu'il montre un six~?}{La chance que les deux soient des six vaut alors $1/6$, puisque le protocole d'information a changé.}
\end{solution}

\begin{solution}{iq:iv1:conditional-probability-and-bayes:8}
La probabilité vaut $2/3$.

Supposons que la bille d'origine soit rouge avec une probabilité de $1/2$. Si elle est rouge, le sac contient deux rouges et le tirage est rouge à coup sûr. Si elle est bleue, le tirage est rouge avec une chance de $1/2$. Donc, par Bayes\index{Bayes par les cotes}, la cote que la bille d'origine était rouge, sachant un tirage rouge, vaut $1 : \tfrac12$, soit $2/3$. Si la bille d'origine était rouge, la bille restée dans le sac est rouge~; si elle était bleue, la bille restante est bleue.

On peut aussi compter les billes directement. Il y a trois rouges équiprobables qui ont pu être la bille tirée~: la rouge ajoutée, la rouge d'origine dans le cas rouge, et la rouge ajoutée dans le cas bleu. Dans deux des trois cas, une bille rouge reste dans le sac.

\iqwrong{Une moitié, en arguant que le tirage a retiré le rouge ajouté, donc que le sac est tel qu'il était. Le tirage rouge est un indice sur la couleur d'origine, car il est plus probable lorsque le sac contenait deux rouges.}
\end{solution}

\begin{solution}{iq:iv1:conditional-probability-and-bayes:9}
Avec un gardien qui répond au hasard, A reste à $1/3$ et C monte à $2/3$. Si le gardien nomme toujours B lorsqu'il le peut, A et C sont chacun à $1/2$.

La vraisemblance des paroles du gardien décide de tout. Si C est gracié, le gardien doit dire B. Si A est gracié, le gardien a le choix, et un gardien au hasard dit B une fois sur deux. Les poids de A et de C valent donc $\tfrac13 \cdot \tfrac12$ et $\tfrac13 \cdot 1$, et la cote vaut $1 : 2$. La chance de A n'a pas bougé, car le gardien aurait pu dire quelque chose de semblable de toute façon, et la nouvelle fait bouger celle de C.

Si le gardien dit toujours B lorsqu'il le peut, alors B est certain lorsque A est gracié, et les poids sont égaux, si bien que chacun a $1/2$. Et un gardien qui nomme C lorsqu'il le peut donne un résultat frappant~: entendre B signifie alors que C est celui qui est gracié.

\iqwrong{«~Deux prisonniers restent, donc une moitié.~» Cela ne tient que pour une règle particulière du gardien. Une probabilité est fixée non seulement par la liste des possibilités restantes, mais aussi par la règle qui a produit le message\index{protocole d'information}.}
\iqwhy{Lorsqu'on livre un fait choisi parmi ce que la personne sait, il faut se demander comment elle l'a choisi. C'est ce qui distingue un rapport filtré d'un échantillon neutre.}
\end{solution}

\begin{solution}{iq:iv1:conditional-probability-and-bayes:10}
Avec la première règle, cela n'a pas d'importance~: chaque porte a une probabilité de $1/2$. Avec la seconde règle, changer de porte gagne à coup sûr après l'ouverture de la porte 3, et avec une chance de $1/2$ après l'ouverture de la porte 2.

Première règle. L'animateur n'a aucune information, et l'ouverture d'une porte vide est elle-même un indice. Si le lot se trouve derrière la porte 1, une porte vide est certaine~; sinon elle a une chance de $1/2$. Les poids valent donc $\tfrac13 : \tfrac23 \cdot \tfrac12$, ce qui est égal.

Seconde règle. Si le lot se trouve derrière la porte 2, l'animateur doit ouvrir la porte 3. S'il se trouve derrière la porte 1, il y a un choix, et l'animateur ouvre toujours la porte 2. Voir la porte 3 s'ouvrir indique donc que le lot se trouve derrière la porte 2~: changer de porte gagne avec certitude. Voir la porte 2 s'ouvrir laisse la porte 1 ou la porte 3, de poids égal, et l'une ou l'autre aurait pu conduire à ce résultat, si bien que la chance vaut $1/2$.

En moyenne sur ce que fait l'animateur, la porte 3 s'ouvre un tiers du temps et la porte 2 s'ouvre les deux tiers du temps, si bien que changer de porte gagne $\tfrac13 \cdot 1 + \tfrac23 \cdot \tfrac12 = \tfrac23$, comme il se doit.

\iqwrong{Donner $2/3$ quelle que soit la règle. La réponse célèbre suppose un animateur qui sait et n'a aucune préférence. Une préférence ne change pas la moyenne, mais elle change ce que l'on peut déduire de ce que l'on a réellement vu.}
\end{solution}

\begin{solution}{iq:iv1:conditional-probability-and-bayes:11}
La probabilité vaut $13/27$.

Compter, en utilisant 14 types d'enfants, un pour chaque sexe et chaque jour. Il y a $14^2 = 196$ familles équiprobables. Celles qui comptent au moins un garçon né un mardi ont le garçon du mardi en premier (14 familles) ou en second (14 familles), en retirant la famille comptée deux fois~: 27. Parmi elles, les deux enfants sont des garçons dans $7 + 7 - 1 = 13$ cas.

Le jour surprenant compte, car une description plus précise rend moins probable que deux enfants différents puissent y correspondre, ce qui fait passer de «~au moins un garçon~» à «~ce garçon-là en particulier~». Sans jour indiqué, la réponse vaut $1/3$. Avec un trait très rare, comme une date d'anniversaire, la réponse vaut $(2 \cdot 365 - 1)/(4 \cdot 365 - 1)$, soit presque $1/2$.

\iqwrong{Répondre que le jour n'est pas pertinent et donner $1/3$. La question ne porte que sur les sexes, mais une famille de deux garçons a presque deux fois plus de chances ($13/49$ contre $7/49$) de compter un garçon né un mardi.}
\iqfollow{La réponse dépend-elle de la manière dont vous avez appris cette information~?}{Oui. Si un parent choisit au hasard un enfant et annonce son sexe et son jour de naissance, la réponse vaut $1/2$. Seule une règle du type «~on annonce s'il a un garçon né un mardi~» donne $13/27$.}
\end{solution}

\begin{solution}{iq:iv1:conditional-probability-and-bayes:12}
La probabilité vaut $2/3$. Ce n'est pas $0.7$, car la loi a priori uniforme compte toujours, comme si l'on avait observé un pile et une face avant de commencer.

La vraisemblance de 7 piles sur 10 lancers, avec la probabilité $p$, est proportionnelle à $p^7(1 - p)^3$, si bien que la loi a posteriori est une loi bêta\index{loi bêta a posteriori} de paramètres $8$ et $4$. Son espérance, qui est la probabilité d'obtenir pile au lancer suivant, vaut $8/12$. En général, la règle de succession de Laplace\index{règle de succession de Laplace} donne $(h + 1)/(n + 2)$.

Sans donnée, elle donne $1/2$, et elle ne donne jamais 0 ni 1 après un nombre fini de lancers, ce que l'on attend d'une règle utilisée pour les événements rares.

\iqwrong{Utiliser la proportion $0.7$. C'est la valeur la plus vraisemblable de $p$, et non son espérance, et la loi a posteriori n'est pas symétrique autour de son maximum. Un a priori plus fort rapprocherait encore davantage la réponse de la moitié.}
\iqwhy{Le rétrécissement vers une loi a priori est intégré à toute estimation de ce type. Un taux de réussite calculé sur peu de transactions doit être ramené vers un taux de base avant d'être pris au sérieux.}
\end{solution}

\begin{solution}{iq:iv1:conditional-probability-and-bayes:13}
La probabilité vaut $35/52 \approx 0.673$.

On met à jour d'abord, puis on prévoit. Les deux piles ont une probabilité $9/16$ avec la pièce biaisée et $1/4$ avec la pièce équilibrée, donc, par les cotes\index{Bayes par les cotes}, la pièce biaisée a des cotes de $9 : 4$, soit une probabilité de $9/13$. La probabilité d'un troisième pile est alors le mélange\index{loi des probabilités totales}
\[
\tfrac{9}{13} \cdot \tfrac34 + \tfrac{4}{13} \cdot \tfrac12 = \tfrac{35}{52}.
\]
Avant tout lancer, ce mélange vaudrait $5/8$. Chaque pile fait donc monter la prévision suivante, parce que l'on apprend sur la pièce.

\iqwrong{Donner $5/8$ au motif que les lancers d'une pièce sont indépendants. Ils le sont conditionnellement à la pièce, mais lorsque la pièce est inconnue, les lancers précédents apportent une information sur elle.}
\iqfollow{Le troisième lancer est-il indépendant des deux premiers~?}{Pas inconditionnellement, et c'est précisément ce qui les rend utiles. Conditionnellement à la pièce, ils sont indépendants.}
\end{solution}

\begin{solution}{iq:iv1:conditional-probability-and-bayes:14}
Après deux tests positifs, la probabilité vaut $9/20 = 45~\%$. Après un troisième test négatif, elle vaut $1/12 \approx 8.3~\%$.

On utilise les cotes et les rapports de vraisemblance\index{rapport de vraisemblance}. Les cotes a priori valent $1 : 99$. Un résultat positif a un rapport de vraisemblance de $0.9/0.1 = 9$, et deux donnent $81$, si bien que les cotes valent $81 : 99 = 9 : 11$, soit $9/20$. Un test négatif a un rapport de $0.1/0.9 = 1/9$, qui ramène les cotes à $1 : 11$, soit $1/12$.

L'ordre des tests n'y change rien. En fait, un positif et un négatif s'annulent, et il reste un seul positif net sur une condition rare.

\iqwrong{Supposer que deux positifs à 90~\% doivent impliquer environ 99~\%. Le taux de base de 1~\% vous coûte plus que ce que deux bons tests peuvent récupérer.}
\iqwhy{Multiplier les rapports de vraisemblance suppose que les tests soient indépendants conditionnellement à la vérité. Répéter le même test sur le même échantillon reproduirait la même erreur deux fois et n'apporterait guère.}
\end{solution}

\begin{solution}{iq:iv1:conditional-probability-and-bayes:15}
Oui, par le paradoxe de Simpson\index{paradoxe de Simpson}. Par exemple, A réussit 9 transactions sur 10 la première année et 30 sur 100 la seconde, tandis que B réussit 80 transactions sur 100, puis 2 sur 10. A est supérieur chaque année (90~\% contre 80~\%, et 30~\% contre 20~\%), pourtant sur les deux années A réussit 39 transactions sur 110 (35~\%) et B en réussit 82 sur 110 (75~\%).

La cause est le mélange. A a surtout traité pendant l'année difficile et B pendant l'année facile, si bien que le chiffre global compare surtout les années. Si chacun avait traité le même nombre de transactions chaque année, l'inversion ne pourrait pas se produire, car les moyennes à poids égaux conservent l'ordre.

Il faut présenter les taux année par année, avec les effectifs de transactions. Si l'on doit donner un seul chiffre, on donne une moyenne sur un mélange commun d'années, car elle compare à conditions égales. Laquelle est la bonne dépend de la question~: l'année est-elle une condition que l'on peut maîtriser, ou un résultat~?

\iqwrong{Répondre «~non~», car un taux plus élevé dans chaque groupe semble forcer un taux plus élevé globalement. Cela n'est vrai qu'avec les mêmes poids sur chaque groupe.}
\iqwhy{L'attribution de performance de portefeuille et la performance selon les régimes se heurtent constamment à ce piège. Une stratégie peut paraître la meilleure globalement uniquement parce qu'elle s'est retenue pendant les mauvaises périodes, ou l'inverse, et il faut savoir laquelle.}
\end{solution}

\begin{solution}{iq:iv1:conditional-probability-and-bayes:16}
La probabilité vaut $16/35 \approx 0.457$, si bien qu'un signal retenu est plus souvent faux que réel.

On compte\index{taux de fausses découvertes} les issues. Cinquante signaux fonctionnent, et 40 d'entre eux sont retenus. Parmi les 950 qui ne fonctionnent pas, $47.5$ sont retenus. Donc $40$ des $87.5$ signaux retenus sont réels.

Il faut d'abord resserrer le test. Au seuil de 1~\%, en gardant la même puissance par simplicité, les fausses alertes tombent à $9.5$ et la part des vrais signaux grimpe à environ 81~\%. En pratique, la puissance baisse aussi, si bien qu'il faudra peut-être davantage de données, mais le taux de base est rarement une grandeur que l'on peut modifier, alors que le seuil l'est. Si aucun signal ne fonctionne, chaque alerte est fausse, et le test ne prévient pas de cela à lui seul.

\iqwrong{Répondre 95~\%, puisque le test est au seuil de 5~\%. Ce chiffre décrit ce qui arrive aux signaux inutiles, et non la probabilité qu'un signal retenu soit utile. La réponse dépend du nombre d'idées testées et de leur qualité.}
\iqwhy{C'est le calcul qui sous-tend tout processus de recherche~: avec un faible taux de bonnes idées, la plupart des résultats significatifs sont faux, sauf si le test est bien plus strict que le seuil habituel de 5~\%.}
\end{solution}

\begin{solution}{iq:iv1:conditional-probability-and-bayes:17}
Le prix juste vaut environ $0.567$, soit 57 centimes.

Chaque pile multiplie les cotes de la pièce à $0.6$ par un facteur $1.5$, et chaque face par $1/1.5$, si bien que seul compte le nombre net de 4 piles\index{rapport de vraisemblance}~: les cotes valent $1.5^4 = 5.0625 : 1$, soit une probabilité de $0.835$ pour la pièce à $0.6$. La probabilité d'un pile au lancer suivant vaut
\[
0.835 \cdot 0.6 + 0.165 \cdot 0.4 \approx 0.567.
\]
Le prix doit se situer entre $0.4$ et $0.6$, quoi que l'on ait observé, car ce sont les deux seules possibilités. Même avec une preuve très forte, le prix s'approche de $0.6$ par en dessous.

\iqwrong{Donner $0.7$, la fréquence observée. Les croyances énoncées n'admettent que les deux pièces, et la fréquence n'intervient que par sa vraisemblance. Les données s'ajusteraient mieux à une troisième pièce, mais on fixe le prix à l'intérieur des croyances énoncées.}
\iqwhy{Les prix se mettent à jour par ce type de calcul. Un teneur de marché qui affiche un cours sur le résultat suivant détient une croyance sur des types, et non une fréquence, et c'est pourquoi la fair value est une moyenne pondérée plutôt que la dernière exécution.}
\end{solution}

\begin{solution}{iq:iv1:conditional-probability-and-bayes:18}
La probabilité vaut $3/4$.

On met à jour selon le contenu de la boîte. Avec $k$ boules rouges, la probabilité de tirer deux rouges sans remise est $\frac{k(k-1)}{12}$, si bien que les poids de $k = 2, 3, 4$ valent $2, 6, 12$, sur 20. La troisième boule est rouge avec la probabilité $(k - 2)/2$, soit $0$, $1/2$ et $1$. La réponse vaut donc $(0 \cdot 2 + \tfrac12 \cdot 6 + 1 \cdot 12)/20 = 3/4$.

Cela concorde avec la règle de succession de Laplace\index{règle de succession de Laplace}, $(2 + 1)/(2 + 2) = 3/4$, et ce n'est pas un hasard. Si le nombre de boules rouges est uniforme sur $0, \dots, N$, la suite des tirages est échangeable, et elle a la même loi que celle de tirages d'une pièce dont le paramètre $p$ est uniforme. La taille de la boîte n'a pas d'importance.

\iqwrong{Deviner une valeur de $k$ et répondre une moitié ou une certitude. C'est l'incertitude sur $k$ qui fait tout le travail.}
\iqfollow{Et après une boule rouge et une boule bleue~?}{Alors $(1 + 1)/(2 + 2) = 1/2$, par la même règle.}
\end{solution}

\begin{solution}{iq:iv1:conditional-probability-and-bayes:19}
La probabilité de culpabilité vaut environ $1/6$.

On part de la population\index{sophisme du procureur}. Avant le test, cette personne a une probabilité $1/(5 \times 10^6)$ d'être la coupable. En dehors de la coupable, environ 5 des 5 millions d'autres personnes correspondent par hasard. Parmi les quelque 6 personnes de la ville qui correspondent, une est coupable, si bien qu'une personne choisie sans autre preuve et qui correspond a une probabilité d'environ $1/6$ d'être coupable.

Par les cotes~: les cotes a priori valent environ $1 : 5\,000\,000$, et le rapport de vraisemblance vaut $10^6$, si bien que les cotes a posteriori valent $1 : 5$.

L'avocat calcule $P(\text{correspondance} \mid \text{innocent})$ et l'énonce comme s'il s'agissait de $P(\text{innocent} \mid \text{correspondance})$. Les deux diffèrent, car presque tout le monde est innocent.

\iqwrong{Traiter le un sur un million comme probabilité d'innocence. D'autres preuves changent pourtant tout~: avec une probabilité a priori d'une moitié tirée d'autres indices, la correspondance rend la culpabilité presque certaine.}
\iqfollow{Et si toute la ville est testée, et qu'une des correspondances est choisie au hasard~?}{Le nombre d'autres correspondances devient alors aléatoire, à peu près de loi de Poisson de moyenne 5, et la probabilité vaut la moyenne de $1/(1 + K)$, soit $(1 - e^{-5})/5 \approx 0.199$. Elle est proche de $1/6$, sans lui être égale, car le protocole\index{protocole d'information} diffère.}
\end{solution}

\begin{solution}{iq:iv1:conditional-probability-and-bayes:20}
La probabilité d'une véritable habileté vaut environ $0.374$, et la probabilité de battre le marché l'année suivante vaut environ $0.575$.

On compare les deux groupes. Les gérants habiles ayant cinq succès d'affilée représentent $0.1 \cdot 0.7^5 \approx 0.0168$ de tous les gérants, et les gérants sans habileté $0.9 \cdot 0.5^5 \approx 0.0281$. Les cotes valent donc $0.0168 : 0.0281$\index{Bayes par les cotes}, et la probabilité vaut $0.374$. La prévision pour l'année suivante est le mélange $0.374 \cdot 0.7 + 0.626 \cdot 0.5 \approx 0.575$.

Cinq succès d'affilée sont plus rares pour un joueur de pile ou face, mais les joueurs de pile ou face sont neuf fois plus nombreux, si bien que la plupart des gérants ayant ce bilan restent sans habileté.

\iqwrong{Prendre le bilan pour une preuve d'habileté, puisque $0.5^5$ ne vaut qu'environ 3~\%. Les gérants sans habileté sont nombreux, et ces 3~\% s'appliquent à tous.}
\iqwhy{Sélectionner des gérants sur leurs rendements passés revient exactement à ce calcul, et il indique de combien un bilan peut faire évoluer l'estimation~: pas rien, mais moins que ne le suggèrent les chiffres bruts.}
\end{solution}

\begin{solution}{iq:iv1:conditional-probability-and-bayes:21}
Enchérir 0, car toute enchère positive a un profit espéré négatif, qui vaut $-b^2/400$. À $2.5\,V$, il faut enchérir 100, pour un profit espéré de 25.

On conditionne sur l'acceptation du propriétaire\index{sélection adverse}\index{malédiction du vainqueur}. Une enchère de $b$ est acceptée exactement lorsque $V \le b$, ce qui a pour probabilité $b/100$, et, sachant cela, $V$ est uniforme sur $[0, b]$ de moyenne $b/2$. On obtient donc en moyenne quelque chose qui vaut $1.5 \cdot b/2 = 0.75b$, et l'on paie $b$. Chaque enchère acceptée perd en moyenne $0.25b$, et le profit espéré vaut $-(b/100)(0.25b) = -b^2/400$.

À un multiple de 2, l'espérance de la valeur vaut exactement l'enchère, et toutes les enchères sont à l'équilibre. Au-delà, chaque enchère est rentable, et une enchère plus élevée est meilleure. Pour $m = 2.5$, le profit vaut $(b/100)(1.25b - b) = b^2/400$, qui vaut 25 pour l'enchère maximale. Une enchère de 40 pour $m = 1.5$ perd 4 en moyenne.

\iqwrong{Enchérir sur la valeur moyenne~: $V$ vaut en moyenne 50, donc elle vaut 75 pour vous, et toute enchère jusqu'à 75 paraît rentable. Cela ignore ce que l'on apprend du vendeur qui dit oui~: l'affaire ne se conclut que lorsque $V$ est faible.}
\iqfollow{Quel doit être le multiple pour qu'une enchère positive ait du sens~?}{Au-delà de 2. Pour un multiple $k$, une enchère de $b$ rapporte en moyenne $\frac{b}{100} \cdot b\bigl(\frac{k}{2} - 1\bigr)$, qui n'est positive que si $k > 2$.}
\iqwhy{Chaque cotation affichée est exécutée sélectivement. Le même conditionnement donne la vraie valeur du prix d'exécution, et c'est pourquoi un prix qui paraît bon avant la transaction peut paraître mauvais juste après.}
\end{solution}

\begin{solution}{iq:iv1:conditional-probability-and-bayes:22}
Les dix premiers tirages sont tous rouges avec la probabilité $1/11$. Sachant que le dixième tirage est rouge, le premier était rouge avec la probabilité $2/3$.

La probabilité que tous les tirages soient rouges se multiplie en $\tfrac12 \cdot \tfrac23 \cdot \tfrac34 \cdots \tfrac{10}{11}$, ce qui se télescope en $1/11$.

L'urne semble se souvenir du passé, mais elle est échangeable\index{echangeabilite@échangeabilité}\index{urne de Pólya}~: toute suite comportant le même nombre de rouges a la même probabilité, car numérateurs et dénominateurs n'apparaissent que dans un ordre différent. Seul compte donc le nombre de rouges, et en fait le nombre de rouges après $n$ tirages est uniforme sur $0, \dots, n$. Donc $P(\text{premier rouge} \mid \text{dixième rouge}) = P(\text{premier rouge} \mid \text{deuxième rouge}) = \tfrac12 \cdot \tfrac23 \big/ \tfrac12 = \tfrac23$. Le tirage ultérieur en dit autant sur le premier que le premier en dit sur le suivant.

\iqwrong{Supposer que l'avenir ne peut pas renseigner le passé, et donner $1/2$. Le conditionnement opère dans les deux sens, car les deux tirages dépendent de la même proportion cachée de rouges.}
\iqfollow{À quoi équivaut cette urne~?}{Une pièce dont la probabilité de pile est uniforme entre 0 et 1, et dont les lancers sont indépendants sachant cette probabilité. C'est pourquoi elle coïncide avec la règle de Laplace.}
\end{solution}

\begin{solution}{iq:iv1:conditional-probability-and-bayes:23}
Les lancers ne sont pas indépendants. Deux piles ont la probabilité $1/3$, et la corrélation vaut $1/3$.

On conditionne sur $p$\index{loi des probabilités totales}~: $P(\text{deux piles}) = E[p^2] = \int_0^1 p^2 \, dp = 1/3$. Chaque lancer pris isolément a la probabilité $E[p] = 1/2$, si bien que l'indépendance demanderait $1/4$. Ils sont indépendants sachant $p$, mais pas globalement~: une pile apporte une information sur $p$.

La covariance vaut $E[p^2] - E[p]^2 = \mathrm{Var}(p) = 1/12$, et chaque lancer a une variance de $1/4$. La corrélation vaut donc $(1/12)/(1/4) = 1/3$. Une pièce dont $p$ est connu n'aurait aucune corrélation, puisque la variance de $p$ y est nulle. En ce sens, la corrélation mesure ce qui est inconnu.

\iqfollow{Et avec davantage de lancers~?}{La même corrélation vaut pour toute paire, puisque les lancers sont échangeables, et la corrélation deux à deux ne s'éteint pas avec la distance.}
\iqwhy{Des rendements qui semblent corrélés entre actifs, ou dans le temps, peuvent tous être des tirages indépendants sachant un état caché, comme un régime de volatilité. Une inconnue commune crée une dépendance sans lien direct.}
\end{solution}

\begin{solution}{iq:iv1:conditional-probability-and-bayes:24}
Sachant $X = 2$, $Y$ a pour espérance 1 et pour variance $3/4$. Et $P(Y > 0 \mid X > 0) = 2/3$.

On écrit $Y = \rho X + \sqrt{1 - \rho^2}\, Z$ avec $Z$ indépendant de $X$. C'est la projection\index{normale conditionnelle} de $Y$ sur $X$. Elle donne une espérance conditionnelle de $\rho x$ et une variance de $1 - \rho^2$, qui ne dépend pas de $x$.

Pour le signe, on utilise l'angle dans le plan de $(X, Z)$\index{probabilité d'orthant}. Les événements $X > 0$ et $Y > 0$ sont des demi-plans, dont les normales forment un angle dont le cosinus vaut $\rho$. La direction de $(X, Z)$ est uniforme, donc
\[
P(X > 0, Y > 0) = \frac14 + \frac{\arcsin \rho}{2\pi},
\]
et $P(Y > 0 \mid X > 0) = \tfrac12 + \arcsin(\rho)/\pi$. Avec $\rho = 0.5$, l'arcsinus vaut $\pi/6$, ce qui donne $2/3$. À titre de vérification, $\rho = 0$ donne $1/2$, et $\rho = 1$ donne 1.

\iqwrong{Répondre $0.75$ en lisant $\rho$ comme une hausse linéaire de la probabilité d'accord des signes. Le lien entre accord des signes et corrélation n'est pas linéaire~: pour $\rho = 0.5$, il vaut $2/3$.}
\iqwhy{Cette même formule relie la justesse du signe et la corrélation d'un signal, si bien qu'un accord de $2/3$ correspond à une corrélation de un demi. De même, une justesse de signe de 55~\% ne correspond qu'à une corrélation d'environ $0.16$.}
\end{solution}

\enlargethispage{-3\baselineskip}%
\begin{solution}{iq:iv1:conditional-probability-and-bayes:25}
La moyenne mise à jour vaut $103.84$ et l'écart-type vaut $2.4$.

Pour des croyances gaussiennes et un bruit gaussien, les précisions s'additionnent\index{pondération par la précision}, et la moyenne est la moyenne pondérée par les précisions. L'a priori a une précision de $1/16$, le signal $1/9$. La nouvelle précision vaut donc $25/144$, soit une variance de $5.76$, ou un écart-type de $2.4$. Le poids du signal vaut $(1/9)/(25/144) = 16/25 = 0.64$, si bien que la moyenne se déplace de $0.64 \cdot 6 = 3.84$ à partir de 100.

Les cas extrêmes se comportent comme on s'y attend~: un signal parfait reçoit tout le poids, et un signal inutile n'en reçoit aucun. Et la variance ne dépend pas de ce que dit le signal.

\iqwrong{Une moyenne simple égale à 103, qui ignore les précisions, ou un écart-type final voisin de 3 ou 4. Combiner les deux donne une incertitude plus faible que chacune des deux.}
\iqfollow{Et après un second signal identique~?}{La précision devient $1/16 + 2/9$, et la moyenne se déplace davantage. En général, la précision s'additionne simplement d'un signal à l'autre, pourvu que les bruits soient indépendants.}
\iqwhy{Cette formule sous-tend le filtre de Kalman et tout rétrécissement d'une prévision bruitée vers un a priori.}
\end{solution}

\begin{solution}{iq:iv1:conditional-probability-and-bayes:26}
Le prix de vente sûr le plus bas est 102, et après cette vente il doit monter à environ $103.85$.

Le prix doit égaler l'espérance de l'action sachant qu'un acheteur se présente\index{sélection adverse}. Si la valeur vaut 110, un achat a la probabilité $0.2 + 0.8 \cdot 0.5 = 0.6$, car tous les traders informés achètent. Si elle vaut 90, un achat a la probabilité $0.4$. Donc $P(110 \mid \text{achat}) = 0.6$, et l'espérance sachant un achat vaut $0.6 \cdot 110 + 0.4 \cdot 90 = 102$. Par le même raisonnement, le cours acheteur vaut 98, et la fourchette de 4 rémunère les traders informés.

Chaque transaction apporte une information. Après un achat, la croyance en 110 vaut $0.6$. Un achat supplémentaire a pour cote\index{Bayes par les cotes} $0.6 \cdot 0.6 : 0.4 \cdot 0.4$, soit $9/13$, et le cours vendeur monte à $90 + 20 \cdot 9/13 \approx 103.85$.

\iqwrong{Coter un prix juste de 100 plus une marge choisie pour dégager un profit. Une cotation à 100 perd en moyenne, car les acheteurs sont plus nombreux quand l'action vaut cher. La fourchette vient du conditionnement sur la transaction, et non d'une commission.}
\iqfollow{Quelles sont les limites~?}{Sans trader informé, le cours vendeur vaut 100. Si tous les traders sont informés, il vaut 110, et à ce prix un acheteur ne gagne rien, si bien qu'aucune transaction ne vaut la peine d'être conclue.}
\end{solution}

\begin{solution}{iq:iv1:conditional-probability-and-bayes:27}
Les probabilités valent $49/58 \approx 0.845$ avec des erreurs indépendantes, $0.7$ si elles sont toujours d'accord, et $11/14 \approx 0.786$ avec les chiffres corrélés.

La hausse et la baisse sont équiprobables, il suffit donc de connaître la probabilité de ce qu'on a observé sachant chaque état du monde. Les deux disent hausse dans un monde où le cours monte exactement lorsque les deux ont raison, et dans un monde où il baisse exactement lorsque les deux ont tort. Donc $P(\text{hausse} \mid \text{les deux disent hausse}) = a/(a + b)$, où $a$ et $b$ sont les probabilités que les deux aient raison et que les deux aient tort\index{preuves corrélées}.

Erreurs indépendantes~: $a = 0.49$ et $b = 0.09$. Toujours d'accord~: $a = 0.7$ et $b = 0.3$, et le second analyste n'ajoute rien. Corrélées~: $a = 0.55$ et $b = 0.15$, ce qui donne $11/14$. Ces nombres sont cohérents~: les autres 30~\% des jours se partagent également entre les deux types de désaccord, si bien que chaque analyste a raison sur $0.55 + 0.15 = 0.70$ des jours, le $0.15$ correspondant aux jours où cet analyste seul a raison.

\iqwrong{Multiplier les rapports de vraisemblance des deux analystes, ce qui donne $(7/3)^2$ dans tous les cas. Cela n'est valable que pour des erreurs indépendantes. Des sources qui se répètent n'ajoutent pas d'élément de preuve.}
\iqwhy{Combiner des signaux en additionnant simplement leur force surestime la confiance si les signaux partagent leurs données. C'est l'une des raisons pour lesquelles un portefeuille de dix signaux semblables vaut bien moins que dix signaux indépendants.}
\end{solution}

\begin{solution}{iq:iv1:conditional-probability-and-bayes:28}
Sachant qu'une main contient au moins un as, la probabilité d'en avoir au moins deux vaut environ $0.370$. Sachant que la main possède l'as de pique, elle vaut environ $0.561$.

On compte les mains. Il existe $\binom{52}{13}$ mains, dont $\binom{48}{13}$ ne contiennent aucun as et $4\binom{48}{12}$ en contiennent exactement un. Donc
\[
P(\ge 2 \mid \ge 1) = \frac{\binom{52}{13} - \binom{48}{13} - 4\binom{48}{12}}{\binom{52}{13} - \binom{48}{13}} \approx 0.370.
\]
Sachant que la main possède l'as de pique, les 12 autres cartes sont tirées parmi les 51 restantes, et n'avoir aucun autre as a pour probabilité $\binom{48}{12}/\binom{51}{12}$. La probabilité d'avoir au moins un autre as vaut donc environ $0.561$.

Les deux événements diffèrent\index{conditionnement sur un élément nommé} parce que «~possède l'as de pique~» est plus précis. Une main de deux as correspond à «~possède l'as de pique~» dans la moitié des cas, alors qu'une main d'un seul as n'y correspond que dans un quart des cas. Nommer la carte favorise donc les mains comptant davantage d'as, exactement comme dans le cas du garçon né un mardi.

\iqwrong{Supposer que la couleur ne peut pas compter, car les as sont symétriques. Chaque couleur a la même probabilité, mais une couleur nommée est moins souvent possédée, et davantage pour les mains comptant moins d'as.}
\iqfollow{Et si un joueur choisit de nommer la couleur de son as lorsqu'il en possède un~?}{Alors la couleur n'apporte aucune information supplémentaire sur le nombre d'as, et $0.370$ tient. C'est le protocole\index{protocole d'information} qui tranche.}
\end{solution}

\begin{solution}{iq:iv1:conditional-probability-and-bayes:29}
Garder la porte choisie à la première ouverture, et changer à la seconde. On gagne avec la probabilité $3/4$.

Il existe quatre stratégies. Garder puis garder donne $1/4$, car les actions du présentateur ne renseignent en rien sur la porte choisie. Garder puis changer gagne chaque fois que le premier choix était faux, puisqu'après deux ouvertures la seule porte fermée qui ne soit pas celle du joueur doit cacher le lot~: $3/4$. Changer puis garder donne $3/8$. Changer dans les deux cas donne $5/8$.

La règle consiste à garder sa porte aussi longtemps que possible, puis à changer au dernier moment. À chaque étape, le présentateur est contraint de fournir une information sur les autres portes, tandis que la probabilité de la porte choisie reste bloquée à sa valeur initiale de $1/4$. Il faut donc garder la porte de plus faible probabilité aussi longtemps que possible, et ne changer que lorsque le changement recueille tout le reste\index{protocole d'information}.

Pour vérifier le $5/8$ obtenu en changeant deux fois~: le premier changement a une probabilité $3/8$ de tenir le lot, le présentateur ouvre une porte supplémentaire, et l'on change pour la porte restante. On gagne sauf si l'on tenait le lot au moment du second changement, ce qui laisse $5/8$. Une énumération exacte de tous les chemins et de tous les choix du présentateur confirme chacun des quatre nombres, et une simulation concorde avec la meilleure.

\iqwrong{Changer aussi souvent que possible, parce que changer une fois était avantageux dans le jeu à trois portes. Chaque changement précoce laisse la porte que le présentateur ne peut pas toucher, c'est-à-dire justement celle dont on veut transférer la probabilité plus tard.}
\end{solution}

\begin{solution}{iq:iv1:conditional-probability-and-bayes:30}
Une famille compte en moyenne 2 enfants, la moitié de tous les enfants sont des garçons, et un enfant rencontré au hasard vient en moyenne d'une famille de 3 enfants.

La taille de famille $L$ suit une loi géométrique sur $1, 2, \dots$ avec $P(L = k) = 2^{-k}$, donc $E[L] = 2$. Chaque famille compte exactement un garçon, donc les garçons représentent un enfant sur deux~: la règle ne modifie pas le ratio des sexes, puisque chaque naissance est un tirage indépendant à pile ou face.

Choisir un enfant n'est pas la même chose que choisir une famille. Une famille de $k$ enfants a $k$ fois plus de chances d'être tirée, il s'agit donc d'un échantillon biaisé par la taille\index{paradoxe d'inspection}\index{echantillonnage@échantillonnage biaisé par la taille}~:
\[
\frac{E[L^2]}{E[L]} = \frac{\mathrm{Var}(L) + E[L]^2}{E[L]} = \frac{2 + 4}{2} = 3.
\]
C'est le même effet que l'intervalle plus long dans lequel tombe une arrivée aléatoire à un arrêt de bus. Si toutes les familles comptaient exactement 2 enfants, les deux moyennes coïncideraient.

\iqwrong{Affirmer que la règle d'arrêt favorise les garçons, ou que la famille d'un enfant pris au hasard compte 2 enfants en moyenne. Une règle d'arrêt ne peut jamais biaiser une pièce équilibrée\index{arrêt optionnel}, mais elle modifie la façon dont on voit les familles lorsqu'on les échantillonne par enfant.}
\iqwhy{Chaque fois que l'on échantillonne par élément plutôt que par groupe, les groupes les plus grands sont surpondérés. C'est pourquoi le fonds moyen d'un investisseur est plus grand que le fonds moyen, et pourquoi la transaction moyenne a lieu à un moment plus actif que la minute moyenne.}
\end{solution}

\begin{solution}{iq:iv1:conditional-probability-and-bayes:31}
La corrélation vaut $-1/2$, alors que les scores sont indépendants dans la population.

Les candidats retenus sont répartis uniformément sur le triangle situé au-dessus de la droite $x + y = 1$. Sur ce triangle, $E[X] = 2/3$, $E[X^2] = 1/2$, donc $\mathrm{Var}(X) = 1/18$. En intégrant, on obtient $E[XY] = 5/12$, donc la covariance vaut $5/12 - 4/9 = -1/36$, et la corrélation vaut $(-1/36)/(1/18) = -1/2$.

La raison en est qu'une personne de faible $X$ ne peut être retenue qu'avec un $Y$ élevé. Cet effet de sélection\index{biais de Berkson}\index{biais de sélection} apparaît avec tout seuil imposé à la somme. Il croît à mesure que le seuil monte jusqu'à 1, puis reste exactement égal à $-1/2$, car au-delà de 1 les candidats retenus forment une copie réduite du même triangle.

Étudier uniquement les candidats retenus peut donc faire croire que les deux compétences s'opposent, ou que l'une d'elles prédit une mauvaise performance. C'est une propriété du filtre, et non des personnes.

\iqwrong{Supposer qu'une sélection préserve l'indépendance. Filtrer selon une fonction des deux variables crée une dépendance là où il n'y en avait pas.}
\iqwhy{Le même piège apparaît lorsqu'une étude n'utilise que les fonds survivants, ou que les transactions exécutées. Tout filtre sur le résultat peut produire une corrélation entre deux variables d'entrée qui n'ont aucun lien entre elles.}
\end{solution}

\begin{solution}{iq:iv1:conditional-probability-and-bayes:32}
Après un rendement de $-4~\%$, la probabilité de crise vaut environ $0.956$. Après un rendement nul, elle tombe à environ $0.017$.

La mise à jour requiert le rapport des densités\index{rapport de vraisemblance} au rendement observé~:
\[
\frac{f_{\text{crise}}(r)}{f_{\text{calme}}(r)} = \frac{1/3 \cdot e^{-r^2/18}}{e^{-r^2/2}} = \frac13 e^{\,4r^2/9}.
\]
Pour $r = 4$, avec des rendements en pourcentage, cela vaut $\tfrac13 e^{7.11} \approx 409$. La cote a priori est de $5 : 95$, donc la cote a posteriori vaut environ $21.5 : 1$, ce qui donne $0.956$. Un seul mouvement peut ainsi accomplir presque tout le travail, puisqu'un mouvement de 4 écarts types est extrêmement rare dans le régime calme.

Pour $r = 0$, le rapport vaut $1/3$, car un rendement proche de zéro est trois fois plus probable dans le régime calme, où la densité est plus pointue. La cote devient $5 : 285$, soit environ $0.017$. Les deux densités sont égales en $|r| \approx 1.57~\%$, et les rendements inférieurs à ce seuil diminuent en fait la croyance en une crise.

\iqwrong{Traiter $-4~\%$ comme «~environ 1.3 écart type sous le régime de crise~», et donc comme une preuve peu forte. La preuve réside dans la comparaison entre deux modèles, et un mouvement banal sous l'un et extrême sous l'autre constitue une preuve très forte.}
\iqwhy{C'est une étape du filtre à la base des modèles à régimes et du comportement d'un GARCH. Les journées calmes devraient réduire l'estimation du risque, et un mouvement important devrait la relever rapidement.}
\end{solution}

\begin{solution}{iq:iv1:conditional-probability-and-bayes:33}
Pour des lois exponentielles, $X$ sachant $X + Y = s$ est uniforme sur $[0, s]$. Pour des normales centrées réduites, elle suit une loi normale de moyenne $s/2$ et de variance $1/2$.

Exponentielles~: la densité conjointe en $(x, s - x)$ vaut $e^{-x} e^{-(s - x)} = e^{-s}$, qui ne dépend pas de $x$. Donc tout partage de $s$ est aussi probable\index{conditionnement sur une somme}~: la moyenne vaut $s/2$ et la variance vaut $s^2/12$.

Normales~: la densité est proportionnelle à $e^{-x^2/2 - (s-x)^2/2}$, qui, en complétant le carré, est une loi normale en $x$ de moyenne $s/2$ et de variance $1/2$. La variance ne dépend pas de $s$ du tout.

Pour toute paire de variables indépendantes et identiquement distribuées, la moyenne de $X$ sachant la somme vaut $s/2$ par symétrie, ce qui ne permet donc pas de savoir quelle loi on a. C'est la dispersion qui les distingue.

\iqwrong{S'attendre à ce que $X$ soit proche de $s/2$ dans tous les cas. Pour des exponentielles, une grande somme est tout aussi probable quel que soit le partage, y compris un grand terme et un tout petit. Pour des normales, une grande somme se répartit.}
\iqwhy{Cette distinction compte pour le risque de queue. Avec des rendements à queues épaisses, un grand mouvement sur deux jours provient le plus souvent d'une journée extrême. Avec des rendements normaux, il provient de deux journées assez fortes. Les conclusions sur la trajectoire dépendent du modèle que l'on retient.}
\end{solution}

\begin{solution}{iq:iv1:conditional-probability-and-bayes:34}
L'erreur consiste à traiter «~l'autre enveloppe contient $2x$ ou $x/2$ avec la même probabilité~» comme vrai pour chaque $x$. Il n'existe pas de loi a priori uniforme sur l'ensemble des montants. Avec une loi a priori propre, cela échoue pour certains montants, et ces montants compensent exactement le gain.

Voici l'exemple. Les six cas équiprobables sont~: on tient 1 ou 2 avec la première paire, 2 ou 4 avec la deuxième, et 4 ou 8 avec la troisième. On tient donc 1 et 8 avec une probabilité $1/6$ chacun, et 2 et 4 avec une probabilité $1/3$ chacun.
\begin{itemize}
\item Si l'on tient 1, l'autre contient 2~: un gain de 1.
\item Si l'on tient 2, l'autre contient 1 ou 4~: une moyenne de $2.5$, soit un gain de $0.5$.
\item Si l'on tient 4, l'autre contient 2 ou 8~: une moyenne de 5, soit un gain de 1.
\item Si l'on tient 8, l'autre contient 4~: une perte de 4.
\end{itemize}
Le gain moyen\index{paradoxe des deux enveloppes} vaut $\tfrac16(1) + \tfrac13(0.5) + \tfrac13(1) + \tfrac16(-4) = 0$. Les deux enveloppes ont en moyenne $3.5$. Changer d'enveloppe gagne donc pour presque tous les montants que l'on pourrait voir, et la grosse perte au sommet compense exactement ces gains.

\iqwrong{«~Puisque cela vaut pour chaque $x$, il faut changer.~» Ce n'est pas vrai pour chaque $x$, comme le montre la valeur la plus élevée. Si l'on tentait de supprimer cette valeur, il faudrait une loi a priori non bornée, et alors l'espérance du montant est infinie, ce qui est précisément là où réside le paradoxe.}
\end{solution}

\begin{solution}{iq:iv1:conditional-probability-and-bayes:35}
La probabilité de l'horaire à 30 minutes vaut $2/3$, et l'on attendra en moyenne encore environ $7.5$ minutes.

Une arrivée aléatoire attend uniformément jusqu'à l'intervalle. Attendre 10 minutes sans bus a une probabilité $5/15 = 1/3$ sous l'horaire à 15 minutes, et $20/30 = 2/3$ sous celui à 30 minutes. La cote est donc de $2 : 1$ en faveur de l'horaire à 30 minutes\index{Bayes par les cotes}.

Sachant l'horaire à 15 minutes et l'absence de bus jusqu'ici, l'attente restante est uniforme sur $[0, 5]$, de moyenne $2.5$. Sachant celui à 30 minutes, elle est uniforme sur $[0, 20]$, de moyenne 10. L'attente restante moyenne vaut donc $\tfrac23 \cdot 10 + \tfrac13 \cdot 2.5 = 7.5$ minutes. Avant d'attendre, l'attente moyenne valait $\tfrac12(7.5) + \tfrac12(15) = 11.25$, donc après 10 minutes l'attente restante vaut $7.5$, et non $1.25$~: avoir attendu a appris que les bus sont moins fréquents qu'on ne l'espérait.

\iqwrong{Utiliser la moyenne des deux horaires sans mise à jour, ou supposer que le bus doit être proche. Avoir attendu est une preuve sur l'horaire, et cela maintient l'attente prévue à un niveau élevé.}
\iqwhy{Toute estimation d'un «~temps jusqu'au prochain~» se comporte ainsi. Si une exécution attendue ou une reprise ne se produit pas, la bonne réaction consiste souvent à allonger l'estimation, et non à croire qu'elle doit être proche.}
\end{solution}

\begin{solution}{iq:iv1:conditional-probability-and-bayes:36}
La probabilité que tout le monde retienne le mauvais type vaut $1/5$, et en moyenne $3.6$ personnes donnent leur avis avant la première dont le signal n'a plus d'importance.

On réduit l'état pas à pas\index{cascade d'information}~: soit $d$ le nombre d'avis antérieurs pour un type moins le nombre pour l'autre. Tant que $|d| < 2$, l'avis de chaque personne révèle son signal. Avec $d = 0$, elle suit son propre signal. Avec $d = 1$, un signal qui va dans le même sens ajoute un avis, et un signal contraire donne une égalité, que la règle tranche en faveur de son signal. Mais en $d = 2$, même un signal contraire laisse un signal net en faveur de la foule. Elle ignore donc son propre signal. Son avis ne révèle rien, la personne suivante fait face au même état, et la foule est bloquée.

Ainsi $d$ est une marche aléatoire, qui se rapproche de la vérité avec la probabilité $p = 2/3$ et s'en éloigne avec $q = 1/3$, absorbée en $\pm 2$. En regroupant les pas par paires\index{ruine du joueur}, la probabilité d'être absorbé à la mauvaise extrémité vaut $q^2/(p^2 + q^2) = 1/5$, et le temps moyen d'absorption vaut $2/(p^2 + q^2) = 18/5$.

Simuler des personnes qui effectuent la mise à jour bayésienne complète à partir des avis publics donne les mêmes nombres.

\iqwrong{Penser qu'une foule finit forcément par avoir raison, parce que les signaux sont bons et que les personnes sont nombreuses. Une fois la cascade amorcée, aucune information nouvelle n'atteint le public, alors même que la quasi-totalité des signaux privés reste inutilisée.}
\iqwhy{Les prix peuvent cesser de révéler l'information lorsque les traders recopient les ordres les uns des autres, et un prix maintenu par une cascade peut se briser à la moindre nouvelle.}
\end{solution}

\begin{solution}{iq:iv1:conditional-probability-and-bayes:37}
Le meilleur gain espéré vaut $\$0.08$. Connaître le type de la pièce donnerait $\$0.60$.

On raisonne par tours, en pariant toujours du côté le plus probable compte tenu de ce que l'on a observé\index{mise à jour séquentielle}. Avec des paris de 1~\$ à la cote 1 contre 1, un pari de probabilité $u$ de gagner rapporte $2u - 1$.

Tour 1~: les deux côtés ont $u = 1/2$, donc le gain est nul. Tour 2~: après un lancer, la croyance dans la pièce favorisée vaut $0.6$, et $u = 0.6 \cdot 0.6 + 0.4 \cdot 0.4 = 0.52$, donc le gain vaut $0.04$. Tour 3~: après HH ou TT, qui a une probabilité totale de $0.52$, la cote vaut $9 : 4$, et $u = 7/13$, donc le pari rapporte $1/13$. Après HT ou TH, la croyance revient à l'égalité, et le pari rapporte 0. Le gain espéré du tour 3 vaut $0.52/13 = 0.04$. Au total, cela fait $0.08$.

En connaissant la pièce, on gagne $0.2$ par tour, soit $0.6$ pour trois tours. Cet edge est pour l'essentiel hors d'atteinte, car il faut d'abord apprendre quelle pièce on tient. Une énumération exacte et une simulation de la règle donnent toutes deux $0.08$.

\iqwrong{Traiter la pièce comme équilibrée en moyenne et répondre zéro. La moyenne a priori est équilibrée, mais l'on apprend au fil du jeu, et parier sur ce que l'on apprend a une valeur positive, bien que faible au début.}
\iqwhy{Les premiers tours ne valent presque rien, sauf comme information. C'est pourquoi un trader peut prendre de petites positions pour apprendre, et pourquoi il faut compter cette valeur lorsqu'on juge si une première transaction a du sens.}
\end{solution}

\begin{solution}{iq:iv1:conditional-probability-and-bayes:38}
Avec 60 sur 100, la probabilité que la pièce soit équilibrée vaut environ $0.523$. Avec $5\,100$ sur $10\,000$, elle vaut environ $0.915$, alors que la p-valeur est à peu près la même. C'est le paradoxe de Lindley\index{paradoxe de Lindley}.

La probabilité d'obtenir 60 piles avec une pièce équilibrée vaut $\binom{100}{60}2^{-100} \approx 0.0108$. Sous la loi a priori uniforme, chaque nombre de 0 à 100 est équiprobable, donc cette probabilité vaut $1/101 \approx 0.0099$. Les données sont donc un peu plus probables sous l'hypothèse d'équilibre, ce qui donne $0.523$.

Avec 10\,000 lancers, la probabilité sous l'hypothèse d'équilibre vaut environ $0.00108$ (la densité normale en $z = 2$ divisée par 50), et la probabilité sous la loi uniforme vaut $0.0001$. Les données favorisent maintenant l'équilibre dans un rapport de 11 pour 1.

Les p-valeurs bilatérales valent $0.057$ et $0.047$. Elles mesurent à quel point le nombre observé est surprenant sous l'hypothèse de la pièce équilibrée. Le facteur de Bayes compare deux explications, et une alternative uniforme répartit sa probabilité sur toutes les valeurs, dont la plupart sont très éloignées de $0.51$. Les données sont donc «~peu probables sous l'équilibre, mais encore moins probables sous l'alternative imprécise~».

La conclusion dépend fortement de l'alternative. Face à l'alternative unique $p = 0.6$, le rapport de vraisemblance vaut environ $7.5$ en faveur de $p = 0.6$, et la probabilité de l'équilibre tombe à environ $0.12$.

\iqwrong{Affirmer qu'une p-valeur de 5~\% signifie une probabilité de 95~\% d'un effet réel, ou une probabilité de 5~\% pour l'hypothèse nulle. Elle ne signifie ni l'un ni l'autre, et la probabilité de l'hypothèse nulle peut même augmenter avec davantage de données pour une même p-valeur.}
\iqwhy{De grands échantillons rendent significatif n'importe quel petit effet. Pour une stratégie disposant d'un échantillon très vaste et d'une statistique de Student limite, demander quelle taille d'effet elle implique, et pas seulement la p-valeur.}
\end{solution}

\begin{solution}{iq:iv1:conditional-probability-and-bayes:39}
Le meilleur affiche environ $1.77$, et son ratio de Sharpe réel espéré vaut environ $0.89$. Sans aucune compétence, le meilleur affiche encore environ $1.25$, et son ratio de Sharpe réel espéré vaut 0.

Chaque valeur mesurée égale le ratio de Sharpe réel plus du bruit, de variance totale $0.25 + 0.25 = 0.5$. Le maximum espéré de 100 normales centrées réduites\index{maximum espéré de normales} vaut environ $2.51$, donc le meilleur affiche $2.51\sqrt{0.5} \approx 1.77$.

Réduisons maintenant. Pour des lois normales, la moyenne de la valeur réelle sachant la mesure vaut $\tfrac{0.25}{0.25 + 0.25} = \tfrac12$ de la mesure\index{rétrécissement}, quel que soit le gestionnaire choisi. Sélectionner le meilleur ne change rien à cette règle, puisqu'elle vaut pour chaque valeur, et le ratio de Sharpe réel espéré du meilleur vaut donc $\tfrac12 \cdot 1.77 \approx 0.89$. La moitié de sa qualité apparente tient à la chance.

Sans compétence, il ne reste que le bruit~: le meilleur affiche $2.51 \cdot 0.5 \approx 1.25$, et la valeur réelle vaut exactement 0.

\iqwrong{Prendre le chiffre du meilleur au pied de la lettre, ou retrancher un ajustement pour le nombre d'essais. La bonne correction dépend de la variation réelle existante, et celle-ci s'estime à partir de la loi a priori, que l'on obtient en mesurant la dispersion des 100 gestionnaires.}
\iqfollow{Comment estimeriez-vous la dispersion de la loi a priori~?}{La variance mesurée entre les gestionnaires est la variance réelle plus la variance du bruit, donc on retranche celle du bruit. Un écart-type mesuré voisin de $0.7$ suggère une dispersion réelle voisine de $0.5$.}
\iqwhy{La performance s'estompe après la sélection pour cette raison, sans aucun changement de compétence. Attendez-vous à ce que la moitié de l'avantage d'un backtest disparaisse à ce niveau de bruit, quel que soit le nombre de stratégies testées~; davantage si le bruit est plus grand par rapport à la dispersion réelle.}
\end{solution}

\begin{solution}{iq:iv1:conditional-probability-and-bayes:40}
Oui. Choisissez un seuil aléatoire selon une loi quelconque qui donne une probabilité positive à tout intervalle, et répondez «~plus grand~» si le nombre vu est inférieur à votre seuil. Pour 10 et 20 avec ce seuil, votre probabilité de réussite vaut environ $0.691$.

Appelons les nombres $a < b$, et soit $T$ votre seuil\index{stratégie du seuil aléatoire}. Il y a trois cas. Si $T < a$, vous répondez toujours «~plus petit~», ce qui est juste exactement lorsque vous avez vu $b$~: une probabilité de $1/2$. Si $T > b$, vous répondez toujours «~plus grand~», ce qui est juste exactement lorsque vous avez vu $a$~: également $1/2$. Si $a < T < b$, vous répondez juste quelle que soit la carte retournée. Votre probabilité de réussite vaut donc
\[
\tfrac12 + \tfrac12\, P(a < T < b),
\]
qui est strictement supérieure à $\tfrac12$ dès que $T$ a une probabilité positive sur tout intervalle. Pour 10 et 20, avec $T$ de loi normale de moyenne 15 et d'écart-type 10, $P(10 < T < 20) \approx 0.383$, donc vous obtenez $0.691$.

La garantie ne peut pas être rendue uniforme. Un adversaire qui connaît votre règle peut placer $a$ et $b$ très proches l'un de l'autre, ou très loin de l'endroit où se situe votre seuil, et l'avantage peut être rendu aussi petit qu'on le veut. Il est positif pour chaque couple, mais il n'admet pas de minorant strictement positif.

\iqwrong{«~Vous ne savez rien de la façon dont les nombres ont été choisis, donc vous ne pouvez pas faire mieux que le hasard.~» C'est vrai pour un seuil fixe, car un adversaire peut placer les deux nombres du même côté de ce seuil. La randomisation supprime précisément cette faiblesse.}
\iqwhy{Ne rien savoir de la taille d'un nombre ne signifie pas que l'observation est inutile. Une valeur relative à une référence, même grossière, renseigne sur l'ordre, et cela procure un petit avantage.}
\end{solution}
